\chapter{Architecture detaillee et diagrammes}

\section{Architecture globale du systeme}

\subsection*{Vue d'ensemble des composants}

L'ecosysteme \textit{Sport Academy Pro} est compose de plusieurs services interconnectes:

\begin{lstlisting}[style=codeblock]
+-------------------------+       REST API       +--------------------------+
|                         | <------------------> |                          |
|  Flutter Mobile App     |                      |  Spring Boot Backend     |
|  (Joueurs, Parents)     |                      |  (football-academy)      |
|                         |                      |                          |
+-------------------------+                      +--------------------------+
         |                                                |
         | WebSocket (Chat)                               | REST API
         v                                                v
+-------------------------+                      +--------------------------+
|                         |                      |                          |
|  Flutter Admin Web      |                      |  MySQL Database          |
|  (Administrateurs)      |                      |                          |
|                         |                      +--------------------------+
+-------------------------+                                |
         |                                                |
         | REST API                                        | JPA/Hibernate
         v                                                v
+-------------------------+                      +--------------------------+
|                         |                      |                          |
|  Django Chatbot Service  |                      |  JMS Broker (Artemis)    |
|  (Assistance Q\&A)      |                      |  (Messaging async)       |
|                         |                      |                          |
+-------------------------+                      +--------------------------+
         |
         | REST API
         v
+-------------------------+
|                         |
|  FastAPI Scouting AI    |
|  (ML: Potentiel,        |
|   Evolution, Churn)     |
|                         |
+-------------------------+
\end{lstlisting}

\section{Architecture microservices}

\subsection*{Separation des responsabilites}

Chaque service a une responsabilite clairement definie:

\begin{table}[H]
\centering
\begin{tabular}{p{0.25\textwidth} p{0.35\textwidth} p{0.3\textwidth}}
\toprule
\textbf{Service} & \textbf{Responsabilite} & \textbf{Technologie} \\
\midrule
football-academy & API metier, authentification, gestion & Spring Boot \\
moez\_project & Interface mobile et web admin & Flutter \\
chatbot & Assistance conversationnelle & Django/ML \\
scouting-ai & Analyse ML et scouting & FastAPI/ML \\
MySQL & Persistence des donnees & MySQL 8.0 \\
Artemis & Messaging asynchrone & JMS \\
\bottomrule
\end{tabular}
\caption{Services et responsabilites}
\end{table}

\section{Flux de donnees principaux}

\subsection*{Flux d'authentification}

\begin{enumerate}[leftmargin=*]
  \item L'utilisateur saisit ses identifiants dans l'app Flutter
  \item Flutter envoie \texttt{POST /api/auth/login} au backend Spring Boot
  \item Spring Boot valide les credentials et genere un JWT
  \item Le JWT est retourne et stocke dans Flutter
  \item Les requetes subsequentes incluent le JWT dans le header \texttt{Authorization: Bearer <token>}
  \item Le \texttt{JwtAuthenticationFilter} valide le token a chaque requete
\end{enumerate}

\subsection*{Flux de messagerie temps reel}

\begin{enumerate}[leftmargin=*]
  \item L'utilisateur ouvre l'ecran de chat dans Flutter
  \item Flutter se connecte au WebSocket endpoint \texttt{/ws}
  \item Flutter s'abonne aux topics STOMP (\texttt{/topic/messages/\{conversationId\}})
  \item Lorsqu'un message est envoye:
    \begin{itemize}[leftmargin=*]
      \item Flutter envoie \texttt{POST /api/chat/messages}
      \item Le backend sauvegarde le message en base
      \item Le backend publie le message sur le topic JMS
      \item Le WebSocket broker diffuse le message aux abonnes
    \end{itemize}
  \item Les clients abonnes recoivent le message en temps reel
\end{enumerate}

\subsection*{Flux de synchronisation IA}

\begin{enumerate}[leftmargin=*]
  \item Le service scouting-ai appelle \texttt{POST /api/v1/data/sync/football-academy}
  \item Il recupere les donnees joueurs depuis \texttt{/football-academy/api/players}
  \item Il recupere les donnees paiements depuis \texttt{/football-academy/api/payments}
  \item Les donnees sont stockees dans la base du service IA
  \item Les modeles ML sont entraines ou mis a jour
  \item Les scores de potentiel, evolution et churn sont calcules
\end{enumerate}

\section{Architecture de securite}

\subsection*{Couches de securite}

\begin{itemize}[leftmargin=*]
  \item \textbf{Authentification}: JWT avec expiration et refresh
  \item \textbf{Autorisation}: RBAC (Role-Based Access Control)
  \item \textbf{Validation}: Validation des entrees et sanitization
  \item \textbf{Encryption}: Mots de passe haches avec bcrypt
  \item \textbf{HTTPS}: Communication securisee (en production)
\end{itemize}

\subsection*{Roles et permissions}

\begin{table}[H]
\centering
\begin{tabular}{p{0.2\textwidth} p{0.7\textwidth}}
\toprule
\textbf{Role} & \textbf{Permissions} \\
\midrule
ADMIN & Acces total, gestion utilisateurs, configuration globale \\
TRAINER & Gestion equipes, planification, saisie statistiques \\
PARENT & Consultation enfants, paiements, planning \\
PLAYER & Consultation profil, planning, communication \\
SCOUTER & Recherche joueurs, comparaison, shortlists (lecture seule) \\
\bottomrule
\end{tabular}
\caption{Roles et permissions}
\end{table}

\section{Architecture de base de donnees}

\subsection*{Schema principal}

La base de donnees MySQL est organisee autour des entites principales:

\begin{itemize}[leftmargin=*]
  \item \textbf{Users}: Utilisateurs du systeme
  \item \textbf{Players}: Profils des joueurs
  \item \textbf{Teams}: Equipes et divisions
  \item \textbf{Activities}: Entrainements, matchs, evenements
  \item \textbf{Payments}: Transactions et cotisations
  \item \textbf{Conversations}: Discussions et messages
  \item \textbf{Notifications}: Alertes et rappels
\end{itemize}

\subsection*{Relations cles}

\begin{lstlisting}[style=codeblock]
User (1) ----< (N) Player
Team (1) ----< (N) Player
Activity (1) ----< (N) Attendance
Player (1) ----< (N) Payment
User (1) ----< (N) Conversation
Conversation (1) ----< (N) Message
\end{lstlisting}

\section{Architecture de deploiement}

\subsection*{Docker Compose}

L'ensemble des services est deploye via Docker Compose:

\begin{lstlisting}[style=codeblock]
version: '3.8'

services:
  mysql:
    image: mysql:8.0
    environment:
      MYSQL_ROOT_PASSWORD: root
      MYSQL_DATABASE: football_academy
    ports:
      - "3306:3306"

  football-academy:
    build: ./football-academy
    ports:
      - "8091:8091"
    depends_on:
      - mysql
      - artemis

  chatbot:
    build: ./chatbot
    ports:
      - "8000:8000"

  scouting-ai:
    build: ./scouting-ai-service
    ports:
      - "8010:8010"
    depends_on:
      - football-academy

  artemis:
    image: apache/activemq-artemis
    ports:
      - "61616:61616"
\end{lstlisting}

\section{Monitoring et observabilite}

\subsection*{Logs et traces}

\begin{itemize}[leftmargin=*]
  \item Logs applicatifs structure (JSON)
  \item Traces des requetes API
  \item Monitoring des performances
  \item Alertes sur erreurs et anomalies
\end{itemize}

\subsection*{Health checks}

Chaque service expose un endpoint de health check:

\begin{itemize}[leftmargin=*]
  \texttt{GET /health} pour football-academy
  \item \texttt{GET /api/health} pour scouting-ai
  \item Verification de la connexion base de donnees
  \item Verification des dependances externes
\end{itemize}